\section{Introduction}\label{sec:intro}

The product of two $n$-bit integers fits in $2n$ bits. Reading the
operands and writing a string of that length takes only linear time.
How much more work does it take to determine the bits of the product?
On a tape the question concerns movement as well as arithmetic. A
Fourier multiplication algorithm repeatedly combines digits stored
far apart, and each tape head moves only one cell at a time.

We construct one machine whose worst-case time per input bit grows
at most as $(1+\log n)^{1/2+o(1)}$. The program, the alphabet and the
number of tapes are fixed once for all lengths. Under a hypothesis
about the time needed to solve satisfiability, the square-root
exponent is optimal.

\begin{theorem}[One multiplier for all lengths]\label{thm:uniform-main}
There is a deterministic multitape Turing machine that computes the
full product of two $n$-bit integers in worst-case time at most
$n(1+\log n)^{1/2+o(1)}$, visiting $O(n)$ tape cells.
\end{theorem}

The assumption for the converse is \emph{variable-count Weak ETH}.
It says that for no fixed $0<\delta<1$ can a deterministic algorithm
decide 3-SAT on $v$ variables in time
$2^{O(v^{1-\delta})}$ times a polynomial in the input length
\cite{SH,WeakSH}. For example, a running time of $2^{O(v/\log v)}$
is compatible with this hypothesis.

\begin{theorem}[Multiplication lower bound]\label{thm:lower}
Assume variable-count Weak ETH. For every deterministic multitape
machine that multiplies exactly there are constants $c>0$ and $n_0$
such that, for each $n\ge n_0$, some pair of $n$-bit operands
requires at least $c\,n\sqrt{\log n}$ steps.
\end{theorem}

\needspace{4\baselineskip}
Both theorems use a fixed finite number of one-dimensional tapes, a
fixed finite alphabet, and one head per tape. Each step reads or
writes at the heads and moves each head by at most one cell. The
operands are nonnegative integers given in binary, high bits first,
with a delimiter between them. The inputs and outputs occupy
prescribed contiguous tape intervals, one bit per cell. Leading
zeros are allowed, and multiplication returns the full $2n$-bit
product. Space counts the cells ever visited, including temporary
work tapes. All logarithms are to base two.

The upper bound in Theorem~\ref{thm:uniform-main} means that the
\emph{same} machine takes $O_\epsilon(n(1+\log n)^{1/2+\epsilon})$
time for every fixed $\epsilon>0$. Whether $O(n\sqrt{\log n})$ time
is attainable remains open. Theorem~\ref{thm:lower} is a statement
at every sufficiently large length, with constants that may depend
on the machine.

\subsection{Division, square root and transposition}

Exact division with remainder and the floor of the square root each
have a single linear-space algorithm with the time bound of
Theorem~\ref{thm:uniform-main}. Newton iteration reduces both tasks
to products at a sequence of increasing precisions. There are also
linear-time reductions in the reverse direction on exact binary inputs.
One division, or three integer square roots, suffices to recover a
product after padding the operands. The padding can hit every large
target length. Thus under Weak ETH even computing the quotient alone
or the floor square root requires $\Omega(n\sqrt{\log n})$ worst-case
time at every sufficiently large $n$ (Corollaries~\ref{cor:near-root-applications}
and~\ref{cor:division-root-lower}).

Matrix transposition connects the two halves of the proof. Store a
square array on a tape one row after another. A scan sees the entries
in each row consecutively, while consecutive entries of a column are
a whole row apart. Transposing the array writes out those columns
consecutively. The multiplier uses fast changes of order of this kind
to bring together the digits of its short Fourier transforms.
Conversely, Harvey and van der Hoeven showed how to transpose a
binary matrix using integer multiplication \cite{hvhtranspose}.

For a matrix of side $N$, let $n=N^2$ denote its bit count. The
reduction and Theorem~\ref{thm:uniform-main} give a transposition
algorithm with time at most $n(1+\log n)^{1/2+o(1)}$. Our conditional
lower bound for this task holds even on a subset of the inputs. For
each fixed $0\le\rho<1-1/\sqrt2$, a machine that returns the correct
transpose by a common deadline on at least a $2^{-\rho n}$ fraction of
the matrices needs a deadline of $\Omega(n\sqrt{\log n})$, at every
sufficiently large side length (Theorem~\ref{thm:transpose}).

\subsection{Earlier bounds}

Karatsuba and Ofman reduced a product to three products of half the
length. Toom and Cook split the inputs into more pieces and used
evaluation and interpolation to recover their product
\cite{karatsubaofman,toom,cook}. The nearly linear algorithms treat
multiplication as convolution. Sch\"onhage and Strassen obtained
$O(n\log n\log\log n)$ tape time using Fourier transforms over
residue rings \cite{schonhagestrassen}. F\"urer reduced the factor
above $n\log n$ to $2^{O(\log^*n)}$ \cite{furer}. Modular arithmetic
and short lattice vectors gave further constructions
\cite{DeKururSahaSaptharishi2013,HarveyHoevenLecerf2016,HarveyHoeven2019Lattice}.
Harvey and van der Hoeven reached $O(n\log n)$ on a multitape machine
using Gaussian resampling to obtain multidimensional Fourier transforms
of power-of-two side lengths \cite{hvh}.

An earlier fixed-tape algorithm combines resampling with a finite
network that shares butterfly operations and address rearrangements.
It gives worst-case time $O(n(1+\log n)^{1-2^{-182}})$
\cite[Theorem~1]{Tape}. Subsequent optimizations of the finite
schedules give larger numerical savings under explicit assumptions
on the cost of their component operations \cite{ChenTapeSchedules}.
Our implementation uses the original stream algorithms and address
changes. For the main convolution we use short cyclic factors
and apply the Fourier operations to many occupied rows at once.
The companion Fourier paper determines the exponent $1/2$ for exact
binary linear circuits, even when one circuit mixes the inputs of
independent transforms \cite[Theorem~\ext{F}{thm:main}]{oaF}. It
also describes earlier savings in operation counts
\cite{AlmanRao,FiniteFourier,ExplicitFourier,Sussman,Kartal}.
To obtain a tape algorithm we must move the rows into position,
compute the coefficients to sufficient precision, and recover every
product digit exactly.

For online multiplication, which must emit digits as the input arrives,
Paterson, Fischer and Meyer obtained an $\Omega(n\log n)$ bound even
on mean time, building on the overlap argument of Cook and Aanderaa
\cite{CA,PFM}. For moving data, Sto{\ss} proved an
$\Omega(N^2\log N)$ lower bound for transposing an $N\times N$ bit
matrix in a model that transports its entries intact \cite{moveonly}.
An ordinary tape machine may read both operands before writing the
product and may encode several input bits together before moving them.

Network coding gives a different route to lower bounds. Adler,
Harvey, Jain, Kleinberg and Lehman connected the undirected
multiple-unicast conjecture to matrix transposition
\cite{AHJ}. Assuming that conjecture, Afshani, Freksen, Kamma and
Larsen proved an $\Omega(n\log n)$ multiplication lower bound for
Boolean circuits with bounded fan-in and fan-out \cite{AFKL}.
Two independent constructions exhibit undirected networks where coding
exceeds the capacity of fractional routing \cite{ZhangLiLi,BravermanHe}. Our lower bound allows arbitrary local coding and charges for
communicating the resulting bits. It assumes variable-count Weak ETH,
which is equivalent to the Stearns--Hunt hypothesis for CNF satisfiability
\cite[Theorem~2.4]{WeakSH}. Subexponential algorithms for 3-SAT
\cite{Anthropic,falseeth} refute the original Exponential Time Hypothesis
\cite{IP} and are compatible with this weaker assumption.


\subsection{How the algorithm saves tape scans}

To multiply the integers, regard their blocks of binary digits as
coefficients of polynomials. Choose the convolution length large
enough that cyclic convolution gives the ordinary coefficient
product. Once the coefficients are known, a scan propagates the
carries. The work in between is to compute Fourier transforms and
place the digits needed together at each stage. A bound on the number
of linear gates in a Fourier circuit does not charge for moving
their inputs across a tape or for representing a complex scalar to
the required precision. The algorithm in Part~\ref{part:algorithm}
accounts for both costs.

\paragraph{Sharing the movement.}
On a power-of-two square array, transposing amounts to exchanging the
row and column fields of the binary address. Divide each field into
$m$ chunks. Swapping the corresponding chunks in turn gives the
desired permutation, but applying each swap separately entails a new
pass over the array. We instead apply combinations of these swaps to
many independent streams. Each swap has eigenvalues $+1$ and $-1$.
In a basis that diagonalizes all the swaps, their product takes its
sign from the parity of the number of minus signs.

The Fourier companion paper gives a way to compute this parity in
bulk \cite[Section~\ext{F}{sec:repair}]{oaF}. Different sign shifts on
the streams give a binomial distribution of counts at each frequency.
On the central interval the parity sign is exactly a polynomial of
degree $O(\sqrt{m\log m})$ in the shifted queries. The exceptional
streams depend on the frequency. A small set of linear combinations
carries enough information into the common basis to recover the
errors there, and further applications of the polynomial correct
the full batch. This gives about $\sqrt{m\log m}$ stages, each linear
in all $m$ swaps.

The tape algorithm must perform such a stage without paying for $m$
separate rearrangements. It groups transfers between streams into
rectangles of source and destination labels. A change of address
layout at a rectangle endpoint can be used for every query there.
In the new layout the shifts needed for each query are controlled by
fields that a tape scan reads earlier. These shifts cost linear time.
The number of smaller interchanges per stream in a stage grows more
slowly than any prescribed positive power of $m$ for a suitable
choice of the finite construction. This is the saving required for
the recursion in Theorem~\ref{thm:interchange}. Section~\ref{sec:permutations}
uses interchange to implement bitwise addition modulo two on selected
address positions and to put records back in order after each transform.

\paragraph{Sharing the arithmetic.}
The construction uses many short, pairwise coprime cyclic factors.
Their product is large enough to hold the digit convolution, and
their padded ranges still fit in an array of linear bit volume.
The Chinese remainder theorem converts the cyclic convolution to a
convolution on this product. Padding each factor afresh for its
Fourier transform could multiply the space requirements across all the
factors. Instead, a short transform borrows occupied coordinates and
restores their contents. This lets the short transforms use the same
sequence of diagonal phases, address permutations and normalized
Walsh transforms (Lemma~\ref{lem:synchronized-short}).

The one-bit Walsh transform combines two values by taking their normalized
sum and difference. On selected address bits it factors into
transforms on smaller sets of bits, with the unselected bits indexing
many independent copies. Lemma~\ref{lem:tape-walsh} shows how to
gather those copies, share the work and return every record to its
specified place on the tape. Our use of the normalized binary
algorithm of the Fourier paper keeps intermediate errors under
control \cite[Proposition~\ext{F}{prop:binary-template}]{oaF}.
With $O(\log n)$ bits per numerical record the total error in the
convolution coefficients is less than one half. Rounding therefore
recovers the exact coefficients before the carry scan. These steps
give a multiplier for every fixed exponent greater than $1/2$
(Theorem~\ref{thm:integer}).

\paragraph{One machine.}
Choosing a new finite construction for each exponent gives a family
of algorithms. Theorem~\ref{thm:uniform-main} requires a single
program with one linear bound on visited space. The finite
constructions can be generated effectively, and each candidate can
test its layout and precision conditions before use. To accommodate a
candidate with a large space constant, the machine tries it on smaller
chunks of the operands and assembles the product from the results.
It runs the candidates with increasing time allowances on a fixed
number of tapes, abandoning attempts that would cross fixed linear
space limits. Every completed attempt gives the correct product.
For each fixed exponent a suitable candidate eventually fits within
the limits and incurs only a constant scheduling overhead
(Section~\ref{sec:single-machine}).

\subsection{Why a fast transposer would contradict Weak ETH}

To prove the lower bound in Part~\ref{part:lower}, we record how the
heads move in a tape run. For a binary $N\times N$ matrix, split time into
epochs of $N$ steps and each tape into blocks of $N$ cells. Given the
block containing a head at the start of an epoch, the head can visit
only that block and its two neighbors. Fix the sequence of starting
blocks for all heads. Each epoch is now a vertex of a network. An
edge passes the contents of a tape block to the next epoch that can
touch it, and neighboring epochs exchange the control state and head
positions. A time bound $T$ gives $O(N+T/N)$ vertices. Even though
the heads move in response to their data, there are only
$2^{O(1+T/N)}$ such coarse traces. Some trace is shared by enough
successful inputs for the partial transposition bound to apply.

That bound prices the wires of the network using the tensor geometry
of the companion papers \cite{oaG,oaP}. A wire carrying $q$ bits
between two sites costs $q$ times their distance. Whatever is
computed locally, independent bits delivered at prescribed pairs of
sites cannot have greater total price than the wires supplying them.
Weak ETH enters when choosing the sites. Uniform expansion for these
distances would give a power saving in the variable exponent of a
3-SAT algorithm (Theorem~\ref{thm:noports}). In the absence of
expansion we find a placement with short network edges but with
typical row and column terminals about one unit apart.

Every correctly transposed bit must appear at its row terminal and
at its column terminal. If the trace succeeds on at least
$2^{-\rho n}$ of the arrays, for $n=N^2$, the successful arrays
contain at least $(1-\rho)n$ bits of joint entropy. Over repeated
copies this yields independent wires at the corresponding rates.
An auxiliary vertex choosing the rows can supply them for at most
$n/\sqrt2+o(n)$ at our placement. For $\rho<1-1/\sqrt2$ a positive
fraction of the demand is still unpaid. This forces the network to
have $\Omega(N\sqrt{\log N})$ vertices (Theorem~\ref{thm:table}),
and hence forces the claimed transposition time.

Finally, given a multiplier at any operand length $m$, the reduction
of Harvey and van der Hoeven transposes a matrix of side
$\lfloor\sqrt m\rfloor$ using a constant number of multiplications at
that very length and $O(m(1+\log\log(2m)))$ other steps. This extra
cost is smaller than the transposition lower bound. The argument
therefore bounds multiplication at every sufficiently large $m$,
without a monotonicity assumption on its running time
\cite[full version, Corollaries~6.2 and~6.3]{hvhtranspose}.
